\section{Team Member Contributions}
\label{sec:contrib}

\begin{table}[H]
\centering\small
\begin{tabularx}{\textwidth}{@{}L{2.35cm} Y@{}}
\toprule
Member & Main completed contributions\\
\midrule
Fu Yitong & Refined the research questions with Wang Junbo and Li Yitong, and took part in revising the research plan after the proposal feedback in November. Designed the first annotation template (issue, claim, stance, reason, source) and its guidelines, compiled a reference corpus of reasons and sources, and ran alignment meetings on unclear annotation cases. Set up the group repository and a rule-based simulation of five intervention strategies. Surveyed research on user controls with Zhao Dingcheng.\\
\addlinespace
Wang Junbo & Refined the research questions and wrote the parts of the project plan on user goals, the exposure report and control design; took part in revising the research plan after the proposal feedback in November. Built the first exposure-report pipeline with Zhao Dingcheng, connecting its stages end to end and generating the report. Designed the annotation template with Fu Yitong, wrote the consistency-check code and annotated a pilot sample of videos. Prepared the presentation slides on interest versus skew and on the proposed goal-to-action design.\\
\addlinespace
Xu Ruicong & Prepared the ShortVideo and KuaiRand data and a 30-user sample of trajectories, and built the initial simulator with Li Yitong. Implemented the transparent recommender model with numerical self-checks; ran the formation, escape (eleven strategies) and persistence simulations with parameter sweeps; computed the detection sample sizes. Drafted the simulation and detection sections and the simulation appendix of the earlier report version, on which \cref{sec:sim,sec:detection,sec:escape} and \cref{app:sim} build.\\
\addlinespace
Zhao Dingcheng & Built the retrieval and extraction prompts of the first pipeline. Defined the four types of opposing narrative and the gender group-label rule (\cref{tab:taxonomy}, \cref{app:prompt}); wrote the labelling prompts, the two-model agreement analysis and the stratified estimators. Designed the first exposure-report mock-up. Reviewed the literature, drafted the comparison of prior work (\cref{tab:related}), and compiled and verified the reference list. Prepared the presentation slides that follow one simulated viewer through the formation, detection and reduction of skew. Drafted the related-work, taxonomy and design sections of the earlier report version, on which \cref{sec:related,sec:data} build.\\
\addlinespace
Li Yitong & Led the redesign of the study after the proposal feedback; derived the definitions, the decomposition of skew and the propositions (\cref{sec:framework}). Ran the data and model work on a GPU server: exposure statistics, the labelling pipeline, and the recommender probing. Divided the work, wrote the task descriptions and reviewed all pull requests. Directed, reviewed and revised this final report (text, figures and layout), which was produced with Claude Code.\\
\bottomrule
\end{tabularx}
\end{table}


\section{Declaration of Generative AI Use}
\label{sec:ai}

We used generative AI at three stages of the project. \Cref{tab:ai} lists the tools and what in this report comes from them; \cref{fig:prompts} gives the prompts that produced that material, translated from Chinese. Each assistant worked in long sessions in which most turns were short instructions to continue or revise; every output was reviewed by the member named and, for shared files, by Li Yitong in a pull request before it was kept.

\begin{table}[H]
\caption{Generative AI used in the project.}
\label{tab:ai}
\centering\small
\begin{tabularx}{\textwidth}{@{}L{2.0cm} L{3.4cm} L{2.4cm} Y@{}}
\toprule
Stage & Tool and model & Used by & Material in this report\\
\midrule
Research plan & Claude Code (Claude Opus 5.5) & Fu Yitong, Wang Junbo & research questions, definitions and propositions (\cref{sec:framework}), choice of domains and case study\\
Implemen\-tation & Claude Code (Claude Opus 5.5) & Li Yitong, Zhao Dingcheng & analysis and labelling code; labelling prompts and narrative taxonomy (\cref{app:prompt}); literature search and reference checks; first drafts of the earlier report\\
 & Codex (GPT-6) & Xu Ruicong & simulation code, parameter sweeps and detection calculations; first drafts of the simulation sections\\
 & Qwen3-30B, DeepSeek V4 Pro and V4 Flash, GPT-6.1 Sol; OneRec-1.7B & Li Yitong, Zhao Dingcheng & research instruments, not writing aids: labelling video descriptions (\cref{sec:data}) and the recommender studied (\cref{sec:formation})\\
This report & Claude Code (Claude Opus 5.5, Claude Fable 5.2) & Li Yitong & text, figures and tables of this report and the scripts that draw them\\
\bottomrule
\end{tabularx}
\end{table}

\begin{figure}[H]
\begin{tcolorbox}[enhanced,colback=panel,colframe=rule,boxrule=0.5pt,arc=2pt,left=8pt,right=8pt,top=5pt,bottom=5pt,fontupper=\small]
\setlength{\parskip}{3pt}
\textbf{Research plan (4--7 November)} \textcolor{mute}{Fu Yitong, Wang Junbo $\to$ Claude Code}\\
``We need to revise the research question. It must be well defined rather than vague, some key concepts need rigorous mathematical definitions, and the main research question must not be scattered; it needs a focus.''\quad
``Sometimes the whole platform shows an overall viewpoint shift [\ldots] Can our current theory identify this kind of shift?''\quad
``I want to frame our topic around university students, choose some related issues [\ldots] and divide all the opposing topics in short videos into a few main categories.''\quad
``I think that works; take gender and marriage as the case study.''

\textbf{Implementation (7--13 November)}\\
\textcolor{mute}{Li Yitong $\to$ Claude Code:} ``Check the progress and push it forward.''\quad ``Continue.''\\
\textcolor{mute}{Zhao Dingcheng $\to$ Claude Code:} ``Start D3.''\quad ``Start D4.''\quad ``Start D5.''\quad ``Start D6.''\quad ``Do Appendix A first.''\quad ``Do \#20, the references.''\\
\textcolor{mute}{Xu Ruicong $\to$ Codex:} ``Continue and finish the earlier tasks for me.''\quad ``Re-submit the earlier ones as pull requests.''

\textbf{This report (14--16 November)} \textcolor{mute}{Li Yitong $\to$ Claude Code}\\
``I hope our report will be a polished, high-quality report, similar to an academic research paper, so I want to write it in \LaTeX{}.''\quad
``Can you make the font of all figures and tables match the font of the text?''\quad
``Yes, just delete it, and tighten the whole report into a fairly pure research question.''\quad
``Revise.''\quad ``Continue.''
\end{tcolorbox}
\vspace{-4pt}
\caption{Prompts that produced material in this report, translated from the Chinese originals. Ellipses mark omitted passages.}
\label{fig:prompts}
\end{figure}
